This study demonstrates how explainable artificial intelligence can be systematically transformed into an operational decision support system for agricultural skill development. Rather than treating machine learning as a purely predictive tool, the results show that model outputs can be structured into interpretable and actionable decision logic, supporting real-world implementation.

A key contribution lies in the integration of segmentation, prediction, and explanation into a unified pipeline. While these components are commonly applied separately in agricultural analytics, their combination enables a direct pathway from data analysis to decision-making. This integrated structure addresses a major limitation in existing applications, where predictive performance is often emphasized without a clear mechanism for translating insights into practice.

The decision rule framework derived from SHAP analysis further extends the role of explainable AI \cite{arrieta2020}. Instead of using SHAP solely for post-hoc interpretation, this study operationalizes feature attribution as the foundation for rule-based decision systems. This transformation allows model outputs to be expressed in an interpretable IF--THEN form, making them accessible to practitioners such as agricultural extension officers.

The findings also highlight the importance of modeling heterogeneity. The segmentation results and cluster-level SHAP analysis reveal that the drivers of skill improvement vary substantially across farmer groups. This suggests that global feature importance alone is insufficient for effective decision-making, and that segmentation-aware models are necessary to capture context-specific dynamics.

In addition, the presence of non-linear relationships, particularly in risk-related variables, illustrates the advantage of machine learning approaches over conventional linear methods \cite{liakos2018}. These patterns indicate that certain factors exert disproportionate influence beyond specific thresholds, a structure that cannot be captured through simple regression models.

From a system design perspective, the proposed Segment--Attribute--Recommend framework addresses critical challenges in deploying AI in agricultural contexts. In particular, the framework is designed for usability, allowing decisions to be implemented without requiring direct interaction with underlying models. This is essential in environments where computational resources and technical expertise are limited.

Despite these contributions, several limitations should be noted. The predictive performance ($R^2 \approx 0.62$) indicates that a portion of outcome variability remains unexplained, likely reflecting unobserved behavioral and contextual factors. In addition, the segmentation structure and threshold patterns are derived from a specific dataset and may require recalibration in different settings. Finally, the analysis is based on observational data, and all findings should be interpreted as associative rather than causal.

Overall, this study suggests that the future of artificial intelligence in agriculture lies not only in improving predictive accuracy, but in developing integrated, interpretable, and deployable decision systems. By embedding explainability within the core logic of decision-making, the proposed framework provides a practical pathway for translating AI capabilities into real-world impact.

Recent advances in agricultural AI research highlight a shift from purely predictive models toward integrated decision systems that combine analytics, interpretability, and real-world implementation \cite{yang2025,zhang2026}. This study aligns with this direction by operationalizing explainable AI into a structured decision support framework.